\paragraph{Conjugate-coordinate class times a fuzzy factor: response of the scale of \(K\) to the frame, the exponent of the \(K\) masses and the scalar channel of the static response (conditional; bounded obstruction).}
This paragraph supplements the preceding one, in which the coupling of massive modes to the scale of \(K\) was recorded as the two alternatives (S1), (S2) and not derived. Keep its setting and notation, with \(V=T\), equal weights and \(q\cdot y=0\) where they are used, and write \(X^i=\chi\Lambda^i\) with fixed Hermitian \(\Lambda^i\), \(\bar a=-\tfrac14\operatorname{tr}[\Lambda^i,\Lambda^j]^2\ge0\), \(\bar n=\operatorname{tr}\Lambda^i\Lambda_i\), \([\Lambda^i,[\Lambda_i,\Lambda_A]]=\mu_A^2\Lambda_A\), \(\gamma^{\mu\nu}=\eta_{ab}E^{a\mu}E^{b\nu}\), \(g=\gamma^{-1}\) and \(\phi=\tfrac12\ln(\rho_M/\sqrt{|g|})\), with \(\rho_M\) taken proportional to the density \(\rho\) of the preceding paragraphs, so that \(\phi\) is the field there up to a constant. A mode \(\psi\Lambda_A\) has the mass \(M^2=\chi^2\mu_A^2\) in the operator \(\gamma^{\mu\nu}\partial_\mu\partial_\nu\) of the preceding paragraph. Define the exponent \(\beta_K\) by \(\delta\ln M^2=2\beta_K\delta\phi=\beta_Kt/2m^2\). On the mass shell the source of a free massive mode is then \(X_\beta=T^{\mathrm m}-\tfrac12\beta_K\eta\operatorname{tr}T^{\mathrm m}\), which couples to \(h_\beta=\bar h-\tfrac12\beta_K\eta t\), the perturbation of \(e^{-2\beta_K\phi}\gamma\); (S1) is \(\beta_K=1\) and (S2) is \(\beta_K=0\).

\emph{Tree level.} For a product the operator on \(\psi\Lambda_A\) is the sum of the operator on \(M\) and \(\chi^2\mu_A^2\), and the tree form is exactly \(\tfrac{\rho_M}{2g^2}\operatorname{tr}(\Lambda_A^2)(\gamma^{\mu\nu}\partial_\mu\psi\partial_\nu\psi+\chi^2\mu_A^2\psi^2)\). So at fixed \(\Lambda^i\) and \(\chi\) the mass is rigid in \(\gamma\), which is (S2), and in the frame \(e^{-2\phi}\gamma\), the effective metric of~\cite{Steinacker2026}, it is \(e^{-2\phi}\chi^2\mu_A^2\), the factor stated in~\cite{MantaSteinacker2025b}. At fixed \(\chi\) a loop carries by power counting one factor \(g^2\sqrt{|g|}/\rho_M=g^2e^{-2\phi}\), a trace over \(M\) against the weight \(\rho_M/g^2\), so corrections to \(M^2\) at \(\ell\) loops have the exponent \(-\ell\). A positive \(\beta_K\) can therefore come only from the response of \(\chi\).

\emph{Scaling.} Under \(X^a\to\sqrt\lambda X^a\), the operators of the content of the preceding paragraphs are exactly \(\lambda H_M+\chi^2H_K\). Since the supertrace of the identity vanishes, the one-loop functional, where its proper-time form converges, depends on \(\lambda,\chi\) only through \(\chi^2/\lambda\) (random product matrices, agreement to twelve digits; the invariance under a common rescaling is used in~\cite{MantaSteinacker2025b}). At one loop alone \(\delta\ln\chi=\tfrac12\delta\ln\lambda\) is a flat direction, and for this homogeneous mode \(\phi=\ln\lambda\) up to a constant, so its exponent is \(\tfrac12\).

\emph{Potential of \(\chi\).} From the factorisation of the preceding paragraphs, with equal weights, \(F=(n+1)^2\) and hats for the unit scale \(\chi=1\), the terms of the action density without derivatives of the wave that depend on \(\chi\) are
\[
\frac{\rho_M}{g^2}V_{\mathrm t}(\chi)+U_4+U_2,\qquad
U_4=-\frac{F}{32\pi^2}\sqrt{|g|}\,\hat U\chi^4,\qquad
U_2=-\frac{F}{32\pi^2}\sqrt{|g|}\Bigl[2(\hat V-\hat T)-\tfrac34\vartheta(\hat V-2\hat T)\Bigr]\chi^2,
\]
with \(\hat U=\int_0^\infty\mathrm ds\,s^{-3}(\hat v^K-\hat\tau^K)(s)\), and \(V=F\hat V\chi^2\), \(T=F\hat T\chi^2\) at equal weights. Here \(V_{\mathrm t}\) contains \(\bar a\chi^4\) and any supplied stabiliser, and \(U_2\) uses \(R/6+E_0=2-\tfrac34\vartheta\) and the index density \(-6\vartheta\) of the background. This has the structure of Equation~(7.48) of~\cite{Steinacker2023}, including the negative sign of the torsion term in the Lorentzian section. The background Feynman gauge is used (Hessian checked by finite differences); its dependence at the stationary point is of higher order by the usual argument, not computed. At \(q\cdot y=q\cdot v=0\) the product \(\bar T\cdot\delta T\) vanishes, so with \(\bar h^{22}=\bar h^{11}\) one has \(\delta\ln(\sqrt{|g|}X)=\Omega/m^2\), \(\Omega=\tfrac12\bar h^{00}-\tfrac53\bar h^{11}\), and \(\Omega=-\tfrac12t\) where \(\bar h^{11}\) is negligible, which holds at large \(\tau\) (below). Write \(U_{\chi\chi}=\partial^2U/\partial(\ln\chi)^2\) at a stationary point. For \(q^2\) below the mass of \(\chi\) defined below and \(V=T\), the first-order response gives in this static channel at large \(\tau\)
\[
\beta_K=\frac{4(2U_4+U_2)}{U_{\chi\chi}},
\]
and for the homogeneous mode \(\beta_K=(8U_4+2U_2)/U_{\chi\chi}\), which is \(\tfrac12\) for \(V_{\mathrm t}=0\). Hence at a local minimum, \(U_{\chi\chi}>0\), for \(U_4\le0\) and \(U_2<0\) one has \(\beta_K<0\) for every \(V_{\mathrm t}\). With \(\tilde U=\rho_M\bar a\chi^4/g^2\) and \(\varepsilon=U_4/\tilde U\): for \(V_{\mathrm t}=\bar a\chi^4\), \(\beta_K=-1/(1+\varepsilon)\); adding a tree cubic term, of the form of the cubic term of Equation~\eqref{eq:Smatrix} read on \(K\) and supplied there, gives \(\beta_K=2(2U_4+U_2)/(2\tilde U+2U_4-U_2)\); adding a tree mass term gives \(\beta_K=(2U_4+U_2)/2(\tilde U+U_4)\). Where the tree stabiliser dominates, \(\beta_K\to0\) from below. The value \(\beta_K=1\) is approached in the last case as \(U_4>0\) comes to dominate \(\tilde U\) and \(U_2\), so it needs \(\hat U<0\), and a tree mass of negative sign for the minimum. In the rotating class of~\cite{MantaSteinacker2025b}, Equation~(3.37) there gives \(\chi\propto\lambda^{-1/6}\) at fixed charge and fixed \(\tau\), the exponent \(-\tfrac16\) for the homogeneous mode; the fluctuation operator about a rotating \(K\) is not treated here.

\emph{Explicit \(K\).} For the fuzzy sphere \(S^2_N\), \(\Lambda^i=J^i\) in three of the six directions, the spectra are closed: the vector \(l(l-1)\), \(l(l+1)+2\), \((l+1)(l+2)\) and the spinor \(l^2\), \((l+1)^2\), without negative mode; \(\bar a=N(N^2-1)/8\), \(\bar n=N(N^2-1)/4\), and \(\chi\) is the only singlet of its symmetry group, so the masses share one exponent at this order. Numerically \(\hat U>0\) with \(\hat U/N^2\to1\), \(\hat T/N^2\to4\) and \(\hat V-\hat T\simeq4\ln N+1.016\) (exact spectra, \(N\le512\)), so \(U_4<0\), \(U_2<0\) and \(\beta_K\le-1\) at the minimum for \(V_{\mathrm t}=\bar a\chi^4\), with \(\beta_K\to-1\) at weak coupling \(|\varepsilon|\ll1\). This has \(\hat U>0\) also for the sampled \(S^2_N\times S^2_{N'}\) and the squashed \(\mathrm{su}(3)\) orbits (a finite list, no general statement). For modes between stacks with constant relative fluxes \(f_1\ge f_2,f_3\ge0\), exactly, \(v-\tau=16s_0\prod_\pm\sinh\tfrac12(f_1\pm f_2\pm f_3)s\), a product over the four sign choices, so for \(f_1<f_2+f_3\), the range without a negative mode, \(\hat U<0\) and \(\hat V<\hat T\), and for \(f_1=f_2+f_3\), \(v\equiv\tau\). So the sign of \(\hat U\) depends on \(K\).

\emph{Scale dependence.} With the tree kinetic coefficient \(Z_\chi=\rho_M\chi^2\bar nm^2/g^2\) of \(\ln\chi\) (the loop part is not computed), \(\ln\chi\) has the mass \(m_\chi^2=U_{\chi\chi}/Z_\chi\) in \(q^2\), and integrating it out replaces \(\beta_K\) by \(\beta_Km_\chi^2/(q^2+m_\chi^2)\), with a contact term \(-2\varrho_0^2m^4/Z_\chi(q^2+m_\chi^2)\). So \(M^2\) is rigid in \(\gamma\) above \(m_\chi\), for every \(K\). For \(V_{\mathrm t}=\bar a\chi^4\), \(m^2m_\chi^2=8(1+\varepsilon)\bar a\chi^2/\bar n\), which is \(4(1+\varepsilon)\chi^2\) on \(S^2_N\). So (S1) is not the coupling computed in this class, and (S2) holds above \(m_\chi\) and approximately where a tree stabiliser dominates.

\emph{Scalar channel of the static response.} For dust and \(X_\beta\), with \(g_2=\vartheta^{-2}[1-\tfrac32c_0\operatorname{arctanh}c_0-\tfrac12(\vartheta-1)]\) the finite part of \(\langle(v^0)^4/(v\cdot v)^2\rangle\), in addition to \(g\) of the preceding paragraph,
\[
\frac{I[X_\beta]}{\varrho_0^2}=\varrho\bigl(\tfrac12\beta_K-g\bigr)^2+g-g_2+\varrho\,(g_2-g^2),
\]
with \(g-g_2\simeq2/\sinh^2\tau\) and \(g_2-g^2\simeq-\tfrac43\sinh^{-4}\tau\), which tends to \(\varrho(1-\tfrac12\beta_K)^2\) as \(\tau\to\infty\) (regulated quadrature agrees). So the response relative to the one-coefficient value is \(4g(1-\beta_K)+\beta_K^2\) at \(\varrho=1\), which reproduces \(1\) and \(4g\) for (S1) and (S2). At \(\varrho=1\), \(\bar h^{00}\), \(\bar h^{11}\) and \(t\) are proportional to \((2-\beta_K)g\), \(\tfrac12\beta_K(1-g)\) and \(\beta_K-2g\), and the ratio of the potentials read in the variable to which the source couples is
\[
\frac{h_\beta^{11}}{h_\beta^{00}}=\frac{\beta_K(1+g-\beta_K)}{4g(1-\beta_K)+\beta_K^2}\ \longrightarrow\ \frac{\beta_K}{2-\beta_K},
\]
while \(\bar h^{00}+\bar h^{11}\) does not depend on the reading. At large \(\tau\) the source therefore acts through \(t\) alone: with \(L=32\pi^2m^6/A\) the reduced form is \(S_E=(q^2/4L)t^2-(1-\tfrac12\beta_K)\varrho_0t\), that is \(\tfrac12(A/4\pi^2)\sqrt{|g|}g^{\mu\nu}\partial_\mu\phi\partial_\nu\phi\) in this static channel, with the positive sign in \(S_E\), so that \(t\) is a physical scalar here and the matter charge is \(1-\tfrac12\beta_K\). With a term \(\tfrac12\mathcal Mt^2\), \(q_*^2=2\mathcal ML\), at \(\varrho=1\) the response is \(-L\varrho_0^2[(\tfrac12\beta_K-g)^2/(q^2+q_*^2)+g(1-g)/q^2]\). A heavy \(t\) therefore leaves only the part with \(g(1-g)\simeq2/\sinh^2\tau\) and ratio \(\tfrac12\) for every \(\beta_K\). A light \(t\) gives \(\beta_K/(2-\beta_K)\), so the value \(1\) at large \(\tau\) needs \(\beta_K=1\), and \(1-\beta_K/(2-\beta_K)=2(1-\beta_K)/(2-\beta_K)\).

\emph{Vacuum term.} At a fibre point the frame perturbation has rank one, so \(\sqrt{|g_n|}=m^{-4}(1+x_n)^{-1}\) with \(x_n\) proportional to \(\alpha\cdot v/m\) exactly, and for a real wave the second-order part of \(U_4\) is \(\tfrac12V_0m^{-6}(\alpha\cdot v)^2\), \(V_0=-F\hat U\chi^4/32\pi^2\). This is the expansion of Equation~(4.27) of~\cite{Steinacker2026}, with its moments, and it enters the form at \(n\) of the preceding paragraph as a shift \(B/A\to B/A-q_b^2/q^2\), \(q_b^2=12\hat U\chi^2/m^2(\hat V+2\hat T)\), which is \(\chi^2/m^2\) on \(S^2_N\) at large \(N\). No term \((v\cdot v)(\alpha\cdot\alpha)\) arises, so on the component of \(\alpha\) orthogonal to \(v\) this term gives no mass. Then \(\varrho\to q^2\varrho/(q^2-\varrho q_b^2)\), and at \(V=T\)
\[
-L\varrho_0^2\Bigl[\frac{(\tfrac12\beta_K-g)^2}{q^2-q_b^2+q_*^2}+\frac{g-g_2}{q^2}+\frac{g_2-g^2}{q^2-q_b^2}\Bigr].
\]
So the forms of the preceding paragraph hold for \(q^2\gg|q_b^2|\). Below that scale, for \(\hat U>0\) the static scalar pole has the tachyonic sign unless \(q_*^2\) exceeds \(q_b^2\), and for \(\hat U<0\) it is a Yukawa pole. That work uses \(V_{01}>0\) from \(S_{\mathrm{vac}}=-\int\rho_M\rho^{-2}V_0\) in Equation~(4.46), and in the passages read does not state the sign of \(V_0\); with the same convention the sampled \(K\) with \(\hat U>0\) have \(V_0<0\), normalisations not compared, so no conclusion on that work is drawn. The tree non-derivative terms, the terms that cancel the tadpole of \(U_4\) at \(q=0\), and any term with the same fibre dependence are not computed, and a long-range form needs them or \(\hat U\to0\).

\emph{The inputs \(V=T\) and equal weights.} \(V\) and \(T\) both scale as \(\chi^2\), so the response of \(\chi\) does not move the condition \(V=T\). On \(S^2_N\), \((V-T)/T\simeq(4\ln N+1.016)/4N^2\to0\) (numerical), so \(V=T\) is approached at large \(N\) for this \(K\) and does not need a tuned \(K\); it holds exactly at \(f_1=f_2+f_3\). At \(V=T\), \(U_2=-(3F/128\pi^2)\vartheta\hat T\sqrt{|g|}\chi^2<0\), so (S1) and \(V=T\) are incompatible at a local minimum for \(\hat U\ge0\): the sign of the torsion term that supplies the minimum is the one that makes \(\beta_K\) negative. On \(S^2_N\) the weights of \(D\) reach \(F\hat V\chi^2\), \(F\hat T\chi^2\) up to a relative correction of order \((n/\chi N)^2\) with a logarithm, with ratios \(0.9994\) at \(N=1024\), \(n=10\), \(\chi=1\), and those of \(B_0,B_1\) carry a common factor from the truncation \(J'\le n\), about \(1-3/2n\) (\(0.8634\) there), the same for \(V_1,T_1,T_2\). So equal weights among the loop levels hold for \(\chi N\gg n\gg1\) on this \(K\). For \(V_{\mathrm t}=\bar a\chi^4\) and \(\vartheta\) near \(1\) this condition and \(|\varepsilon|\ll1\) leave the window \(F/4\pi^2N\ll\rho_M/g^2\sqrt{|g|}\ll3NF/8\pi^2n^2\), as a scale estimate. The hypothesis on the external harmonics above \(J=1\) is not tested.

\emph{Consequence.} Within these hypotheses a static, late-time response with ratio \(1\), read in the variable to which the source couples, requires a light \(t\) and \(\beta_K=1\), and the computed exponent is negative for \(\hat U\ge0\). With a light \(t\) and \(\beta_K<0\) the deviation of this ratio from \(1\) exceeds \(1\), and lies in \((1,\tfrac43]\) for \(\beta_K\in[-1,0)\), to be set against the value \((2.1\pm2.3)\times10^{-5}\) of \(\gamma-1\) in~\cite{Bertotti2003}. This is not a Cassini bound on the emergent metric, since solar-system scales lie where \(q^2\ll|q_b^2|\) and the non-derivative terms are not computed. The route to \(\beta_K=1\) that remains here is \(\hat U<0\), a tree mass and a loop quartic above the tree quartic, for which \(q_b^2<0\) and the vacuum term must then be cancelled at the minimum without cancelling its derivative in \(\chi\).

These statements are for the hypotheses of the preceding paragraph, the product background without mixed flux, one loop, the formal continuation and the frozen symbol, for scales of \(q\) as stated; the first-order response in the static channel at \(q\cdot y=0\), with \(V=T\) where the index part alone enters \(U_2\) and with \(W_R\propto\hat V-\hat T\) neglected; the fibre-constant background for \(U_2\); and a mass of the matter modes that comes from \(\chi\) alone, so that masses from gauge dynamics, which see \(g^2e^{-2\phi}\), and the non-singlet moduli of \(K\) are not treated. The sign statement and the exponents are exact for these forms, the values on explicit \(K\) are exact spectra with numerical fits, and the window is a scale estimate. They do not construct the \(K\) selected by the theory, determine \(\mathcal M\) or the background value of \(\phi\), compute the loop kinetic term of \(\chi\), the rotating class, higher loops or the fibre Laplacian, identify \(\bar h\), \(h_\beta\) or \(t\) with quantities of the emergent metric of Chapter~\ref{ch:vierbein-emergence}, discharge Assumption~\ref{ass:gap-psi}, or bear on Assumption~\ref{ass:lorentzian}. Problems~C, F and~N remain open and PROGRAMME-CONDITIONAL.

